\section{Neural Field: cómo un mismo Whole genera distintas Parts}

En los últimos dos minutos se resolvió una cuestión que es fácil pasar por alto, pero que determina si el top-down path es coherente: si la face-level island es el mismo vector en todas las posiciones y todas las columnas comparten la misma top-down network, ¿por qué en unas posiciones la salida es nose y en otras mouth? La respuesta es que la top-down function también recibe la target location.

\subsection{Location-conditioned top-down prediction}

Sea $\mathbf e_w$ el whole embedding y $x$ la posición que se debe predecir; entonces el shared decoder es la función

\[
\hat{\mathbf e}_{\mathrm{part}}(x)=f_{\mathrm{td}}\!\left(\mathbf e_w,x\right).
\]

Aquí, $\mathbf e_w$ codifica a la vez la whole identity y la pose, $x$ especifica la posición de la imagen que se consulta, y la salida es el lower-level identity-pose embedding que debería aparecer en esa posición. El mismo face vector, combinado con distintos valores de $x$, puede generar nose, mouth u otra part prediction; esta es la forma básica de un coordinate-conditioned neural field.

\teachervoice{Hinton plantea primero una contradicción aparente: las red arrows y las green arrows son completamente distintas, pero los face vectors del nivel superior son idénticos; si los top-down weights también se comparten, ¿cómo se obtienen resultados distintos? Su solución es introducir también la query location como entrada. Esta aclaración convierte «pesos compartidos» de un lema en una firma de función.}

\begin{knowledgebox}{Qué resuelve aquí el neural field}
\begin{itemize}
  \item Entrada: el whole identity/pose embedding más la target coordinate;
  \item Salida: la lower-level prediction en esa coordinate;
  \item Valor: con un mismo decoder compartido, se producen partes distintas en posiciones distintas;
  \item No equivale a: la clase no exige usar el volume rendering, el ray sampling ni la photometric loss de NeRF.
\end{itemize}
\end{knowledgebox}

\subsection{Relación con los coordinate-conditioned models actuales}

NeRF, las implicit neural representations y algunos world model decoders usan «global/local latent + coordinate/query» para producir salidas específicas de cada posición. Lo distintivo de GLOM es que el decoder no se limita a renderizar píxeles a partir de un único global latent, sino que está integrado en recurrent islands de varios niveles: la top-down prediction debe converger junto con la bottom-up evidence, el lateral agreement y el temporal state.

\begin{warningbox}{Una Connection no es una lineage claim}
Relacionar GLOM con los object-centric decoders, slot models y world models actuales es una interpretación de ingeniería; no puede afirmarse que estos últimos «implementaron GLOM», salvo que de verdad cumplan condiciones más fuertes, como la island representation, la multi-level recurrence, el four-source consensus y la coordinate-aware part-whole prediction.
\end{warningbox}

\subsection{Resumen de la sección}

La Location condition resuelve la contradicción expresiva de la shared top-down network: una misma whole representation puede emitir distintas part predictions según la coordenada. Es también el componente de todo el diseño más cercano a una técnica moderna consolidada, pero la dificultad de GLOM sigue siendo cómo entrenar ese decoder de forma estable y conjunta con las islands dinámicas.
